%% kapa: exam=1 round=23 year=2012 period=1 form=A subject=economics subject_ko=경제원론 seq=34 number=74
%% source: Q-Net 감정평가사 기출문제 (cst00309) artlSeq=5052578 file=raw/1008/5052578/2012년 제23회 감정평가사 제1차시험 기출문제(1교시 A형).pdf
%% conversion: kordoc 4.24.51 -> tools/gen1.py (auto); review=proofread-sonnet; flags=tikz
\begin{question}{74}
다음 그림은 A국의 인플레이션율과 실업률 사이의 단기적 상충관계를 나타내는 필립스곡선이다. 이 관계에 근거하여 단기적으로 실업률을 낮추기 위한 정부의 정책 방향으로 옳은 것은? (단, 세로축은 인플레이션율, 가로축은 실업률이고, 단위는 \%이다.)
\begin{center}
\begin{tikzpicture}[x=0.1mm,y=-0.1mm,font=\scriptsize,line width=0.9pt]
\draw (65,5) -- (65,440) -- (620,440);
\draw[line width=1.4pt] plot[smooth] coordinates {(108,73) (120,150) (160,235) (220,310) (300,365) (400,395) (520,409)};
\node[align=center,font=\tiny] at (20,150) {인\\플\\레\\이\\션\\율};
\node at (330,478) {실업률};
\node at (47,458) {0};
\end{tikzpicture}
\end{center}
\begin{choices}
  \item 정부지출을 감소시킨다.
  \item 소득세를 인하한다.
  \item 통화량을 감소시킨다.
  \item 기준금리를 인상한다.
  \item 법인세를 인상한다.
\end{choices}
\end{question}
\answer{2}
